\documentclass[11pt]{article}
% Standalone preamble for web-only supplementary notes.
\usepackage[T1]{fontenc}
\usepackage{amsmath,amssymb,amsthm}
\usepackage{hyperref}
\usepackage{cleveref}
\DeclareMathOperator{\Map}{Map}
\DeclareMathOperator{\Fun}{Fun}
\newcommand{\iso}{\xrightarrow{\sim}}
\newtheorem{lemma}{Supplementary Lemma}
\newtheorem{corollary}[lemma]{Supplementary Corollary}
\crefname{lemma}{Supplementary Lemma}{Supplementary Lemmas}
\Crefname{lemma}{Supplementary Lemma}{Supplementary Lemmas}
\crefname{corollary}{Supplementary Corollary}{Supplementary Corollaries}
\Crefname{corollary}{Supplementary Corollary}{Supplementary Corollaries}

\begin{document}

\section*{Putting identity identifications into the inputs}

The \href{coherences.html\#rmk:Isomorphisms_Of_Natural_Isomorphisms}{remark
following horizontal compatibility} says that the individual whiskering
laws follow by taking some inputs to be identity identifications.
There are endpoint comparisons in this substitution: inserting an identity
identification produces whiskering, while inserting an identity functor
changes the endpoints by a unitor. We display the resulting laws here.

We use the structures and coherences of Sections 1.1--1.2.
Write $\tau\cdot\sigma$ for vertical composition, $u\star\sigma$ and
$\sigma\star k$ for whiskering, and $\tau\ast\sigma$ for horizontal
composition. The right-hand identification in a vertical composite is applied
first. These operations also apply to families with a common parameter category
$X$, namely functors from $X$ into fixed identification animae.
For an absolute identification $\delta$, write $\overline\delta$ for its
constant family on $X$.

\begin{lemma}
\label[lemma]{lem:Supplementary_Horizontal_Identity_Inputs}
For $f,f'\colon C\to D$, $g,g'\colon D\to E$, and families
$\sigma\colon X\to(f=f')$ and $\tau\colon X\to(g=g')$, there are
identifications
\[
 \overline{\mathrm{id}_g}\ast\sigma=g\star\sigma,
 \qquad
 \tau\ast\overline{\mathrm{id}_f}=\tau\star f.
\]
In particular,
\[
 \overline{\mathrm{id}_g}\ast\overline{\mathrm{id}_f}
   =\overline{\mathrm{id}_{g\circ f}}.
\]
\end{lemma}
\begin{proof}[Proof idea]
By definition,
\[
 \tau\ast\sigma=(\tau\star f')\cdot(g\star\sigma).
\]
If $\tau$ is the constant identity family, preservation of identities by
prewhiskering identifies the first factor with the constant identity of
$g\circ f'$. Apply the left vertical unit law. For the second formula,
use preservation of identities by postwhiskering and the right vertical unit
law. Constant families are compared with their whiskerings using the
associator. Combining the first formula with preservation of identities by
postwhiskering gives the last formula.
\end{proof}

\section*{The five endpoint-compatible laws}

We use the book's names for the structural identifications:
\[
 \lambda_f\colon\mathrm{id}\circ f=f,\qquad
 \rho_f\colon f\circ\mathrm{id}=f,\qquad
 \alpha_{f,g,h}\colon(h\circ g)\circ f=h\circ(g\circ f).
\]

\begin{lemma}
\label[lemma]{lem:Supplementary_Five_Whiskering_Laws}
Let $f,g\colon C\to D$ and put $X=(f=g)$, with universal family
$\sigma=\mathrm{id}_{(f=g)}$. The unit comparisons give
\[
 \overline{\lambda_g}\cdot(\mathrm{id}_D\star\sigma)
     =\sigma\cdot\overline{\lambda_f},
 \qquad
 \overline{\rho_g}\cdot(\sigma\star\mathrm{id}_C)
     =\sigma\cdot\overline{\rho_f}.
\]
For $u\colon D\to E$ and $v\colon E\to F$, there is an identification
\[
 \overline{\alpha_{g,u,v}}\cdot((v\circ u)\star\sigma)
   =\bigl(v\star(u\star\sigma)\bigr)\cdot\overline{\alpha_{f,u,v}}.
\]
For $k\colon B\to C$ and $l\colon A\to B$, there is an identification
\[
 \overline{\alpha_{l,k,g}}\cdot((\sigma\star k)\star l)
   =(\sigma\star(k\circ l))\cdot\overline{\alpha_{l,k,f}}.
\]
Finally, for $k\colon B\to C$ and $u\colon D\to E$, there is an identification
\[
 \overline{\alpha_{k,g,u}}\cdot((u\star\sigma)\star k)
   =(u\star(\sigma\star k))\cdot\overline{\alpha_{k,f,u}}.
\]
All five are identifications between functors on $(f=g)$.
\end{lemma}
\begin{proof}[Proof idea]
For the unit formulas, start with the horizontal unit comparisons and use
\Cref{lem:Supplementary_Horizontal_Identity_Inputs} to replace horizontal
composition with an identity identification by whiskering.

For the postcomposition formula, substitute
$\overline{\mathrm{id}_v},\overline{\mathrm{id}_u},\sigma$ into horizontal
associativity. The two parenthesizations become whiskering by $v\circ u$
and successive whiskering by $u$ and $v$. The two vertical sides remain
$\overline{\alpha_{f,u,v}}$ and $\overline{\alpha_{g,u,v}}$.
The other two formulas use, respectively,
$\sigma,\overline{\mathrm{id}_k},\overline{\mathrm{id}_l}$ and
$\overline{\mathrm{id}_u},\sigma,\overline{\mathrm{id}_k}$.
In each case the horizontal identity comparisons give the displayed boundary.
\end{proof}

Precomposing these universal comparisons with a family
$R\to(f=g)$ gives the corresponding formulas for that family. The endpoint
normalizations are the ones explained in
\href{supplementary-specializing-coherence-laws.html}{Specializing coherence laws}.
This derivation uses horizontal units and associativity together with the
vertical laws and preservation of identity identifications by whiskering.
It does not require interchange.

\section*{The middle interchange in a composite}

The book also observes that horizontal composition preserves vertical
composition. The following calculation makes explicit the single interchange
needed for absolute identifications.

\begin{lemma}
\label[lemma]{lem:Supplementary_Horizontal_Preserves_Vertical}
Let $f_0,f_1,f_2\colon C\to D$ and $g_0,g_1,g_2\colon D\to E$.
For absolute identifications $\sigma_i\colon f_{i-1}=f_i$ and
$\tau_i\colon g_{i-1}=g_i$, for $i=1,2$, there is an identification
\[
 (\tau_2\cdot\tau_1)\ast(\sigma_2\cdot\sigma_1)
   =(\tau_2\ast\sigma_2)\cdot(\tau_1\ast\sigma_1).
\]
Horizontal composition also preserves identity identifications:
$\mathrm{id}_g\ast\mathrm{id}_f=\mathrm{id}_{g\circ f}$.
\end{lemma}
\begin{proof}[Proof idea]
Expand the horizontal composite on the left and use preservation of vertical
composition by each whiskering. This gives the four factors, in their order
of vertical composition,
\[
 (\tau_2\star f_2),\quad(\tau_1\star f_2),\quad
 (g_0\star\sigma_2),\quad(g_0\star\sigma_1).
\]
The fixed-outer interchange comparison for $\tau_1$, specialized at
$\sigma_2$, identifies the middle composite with
\[
 (\tau_1\star f_2)\cdot(g_0\star\sigma_2)
   =(g_1\star\sigma_2)\cdot(\tau_1\star f_1).
\]
Reassociate the four factors before and after this replacement. The result
is exactly the right-hand horizontal composite. Preservation of identities
follows by expanding its two whiskerings and then applying a vertical unit law.
\end{proof}

Here the inputs $\tau_i,\sigma_i$ are absolute identifications. The proof
uses the stated fixed-input interchange axiom. A comparison over a parameter
category in which both interchange inputs vary is a stronger assertion;
the formalization treats that assertion separately under an additional joint
interchange hypothesis. The displayed result does not use that hypothesis.

\end{document}
